\section{Preface}
\lipsum[1]

Testing some equations: let $\vec{x},\vec{y}\in\Rs[n]$ and $\alpha,\beta\in\Rs$. Then from the fact that $\Rs[n]$ is a vector space, we get
\begin{equation}
	\vec{c} = \alpha\vec{x} + \beta\vec{y} \in \Rs[n].
	\label{eq:vector_space_closure}
\end{equation}

However, if $\vec{a}\in\Rs[m]$, where $m\neq n$, then $\alpha\vec{x}\in\Rs[m]$, and therefore $\vec{c}$ is not defined (i.e. \autoref{eq:vector_space_closure} is irrelevant).

Now, this is a test of the \textit{aligned} environment, which allows for multiple (correctly-aligned) lines of equations and having the label only once (instead of per line, as in the \textit{align} environment):
\begin{equation}
	\begin{aligned}
		\sin(2x) & = \frac{1}{2\iu}\left(\Eu{2x\iu}-\Eu{-2x\iu}\right)                                                                             \\
		         & = \frac{1}{2\iu}\left(\left[\Eu{x\iu}\right]^{2}-\left[\Eu{-x\iu}\right]^{2}\right)                                             \\
		         & = \underbrace{\frac{1}{2\iu}\left(\Eu{x\iu}-\Eu{x\iu}\right)}_{\sin(x)}\underbrace{\left(\Eu{x\iu}+\Eu{x\iu}\right)}_{2\cos(x)} \\
		         & = 2\sin(x)\cos(x).
	\end{aligned}
	\label{eq:aligned_env}
\end{equation}

\begin{note}{Test note}{test_note}
	This is a test note. I will try to also refernce it later.
\end{note}

Now, I wish to refernce \autoref{note:test_note}. Did it work?

Also, let us look at an example.

\begin{example}{Some test example}{test_example}
	If $x=\rowvec{1;2;-3}$ and $y=\rowvec{5;-1}$, then $n=3$ and $m=2$. Clearly, $n\neq m$ - and so $\alpha\vec{x}\in\Rs[3]$, while $\beta\vec{y}\in\Rs[2]$. This means that $\alpha\vec{x}+\beta\vec{y}$ is not defined.
\end{example}

\dots let's check if referencing an example works correctly: can you see \autoref{example:test_example}?

I believe it is now time for some figures! See \autoref{fig:test_fig_1}.

\lipsum[1]

\begin{wrapfigure}{l}{75mm}
	\begin{center}
		\begin{tikzpicture}
			\coordinate (o) at (0,0);
			\draw[vector, xred] (0,0) -- (3,1) node [pos=1.1] (a) {$\vec{w}_{1}$};
			\draw[vector, xblue] (0,0) -- (2,3) node [pos=1.1] (b){$\vec{w}_{2}$};
			\pic[draw, xpurple, -stealth, thick, angle radius=2cm, angle eccentricity=1.1, pic text=$\varphi$]{angle=a--o--b};
			\draw[vector, xgreen] (0,0) -- (-1,2) node [pos=1.1] (c) {$\vec{x}_{1}$};
			\draw[vector, xorange] (0,0) -- (-3,2) node [pos=1.1] (d){$\vec{x}_{2}$};
			\pic[draw, black!75, -stealth, thick, angle radius=1.25cm, angle eccentricity=1.2, pic text=$\theta$]{angle=c--o--d};
		\end{tikzpicture}
	\end{center}
	\caption{This is the first test figure. Does it look ok?}
	\label{fig:test_fig_1}
\end{wrapfigure}

\lipsum[1]

\begin{wrapfigure}{c}{90mm}
	\begin{center}
		\begin{tikzpicture}
			\begin{axis}[
					width=10cm,
					xyplane,
					xmin=-2, xmax=2,
					ymin=-2, ymax=2,
					xtick={-2,...,2},
					ytick={-2,...,2},
					xticklabels={,$-1$,,$1$,},
					minor tick num=4,
					samples=150,
				]
				\addplot[name path=sqrd, very thick, xred] {\x^2-\x-1};
				\addplot[name path=cubd, very thick, xblue] {\x^3-\x};
				\addplot[thick, fill=xgreen, fill opacity=0.4] fill between[of=sqrd and cubd, soft clip={domain=0:1}];
				\coordinate (a) at (-0.755,0.325);
				\addplot [only marks, mark=*, mark options={fill=white, scale=1.5}] table {
						-0.755 0.325
					};
				\node[above left of=a, xshift=-20pt, align=right] (a_txt) {intersection};
				\draw[-stealth] (a_txt.south) to [out=270, in=180] (a);
			\end{axis}
		\end{tikzpicture}
	\end{center}
	\caption{And this is another test figure! I hope it works.}
	\label{fig:test_fig_2}
\end{wrapfigure}

\lipsum{1}
